%=============================================================================!
% 18.F  SOLUTIONS TO THE EXERCISES                                            !
%=============================================================================!
\unnumbered{Chapter~\ref{ch:learning}: Learning as optimisation}
\label{sec:lo-solutions}

\solutionto{ex:lo-least}With $S_{xx} = \sum_ix_i^2$, $S_{xF} =
\sum_ix_iF_i$ and $S_{FF} = \sum_iF_i^2$, the expanded loss is $L = S_{FF}
- 2kS_{xF} + k^2S_{xx}$, and at $\hat k = S_{xF}/S_{xx}$
\begin{equation*}
   L(\hat k) = S_{FF} - 2\frac{S_{xF}^2}{S_{xx}} + \frac{S_{xF}^2}{S_{xx}^2}
   S_{xx} = S_{FF} - \frac{S_{xF}^2}{S_{xx}} .
\end{equation*}

\solutionto{ex:lo-error}The stretches in metres are 0.01 to 0.10, and
the sum of their squares is $385\times10^{-4} =
\SI{0.0385}{\metre\squared}$, so the error is $0.05/\sqrt{0.0385} =
\SI{0.255}{\newton\per\metre}$.

\solutionto{ex:lo-intercept}The columns of $X$ are 1 and $x_i$, so
\begin{equation*}
   X^{\mathsf T}X = \begin{pmatrix} n & \sum_ix_i\\ \sum_ix_i &
   \sum_ix_i^2\end{pmatrix} , \qquad
   X^{\mathsf T}\vb{y} = \begin{pmatrix}\sum_iy_i\\ \sum_ix_iy_i
   \end{pmatrix} .
\end{equation*}
The first equation, $nc + s\sum_ix_i = \sum_iy_i$, divided by $n$ gives
$c = \bar y - s\bar x$. Substituting into the second, $(\bar y - s\bar
x)\sum_ix_i + s\sum_ix_i^2 = \sum_ix_iy_i$, and with $\sum_ix_i = n\bar x$,
\begin{equation*}
   s\Bigl(\sum_ix_i^2 - n\bar x^2\Bigr) = \sum_ix_iy_i - n\bar x\bar y .
\end{equation*}
Expanding $\sum_i(x_i - \bar x)^2 = \sum_ix_i^2 - 2\bar x\,n\bar x + n\bar
x^2 = \sum_ix_i^2 - n\bar x^2$, and likewise $\sum_i(x_i - \bar x)(y_i -
\bar y) = \sum_ix_iy_i - n\bar x\bar y$, gives the slope of
\cref{sec:cv-drift}.

\solutionto{ex:lo-stiffness}With $U_{\mathrm M} = D(1 -
\mathrm{e}^{-ax})^2$, the chain rule gives $F = -2Da\,\mathrm{e}^{-ax}(1 -
\mathrm{e}^{-ax}) = -2Da(\mathrm{e}^{-ax} - \mathrm{e}^{-2ax})$, and
\begin{equation*}
   -\frac{\dd F}{\dd x} = 2Da(-a\mathrm{e}^{-ax} + 2a\mathrm{e}^{-2ax})
   = 2Da^2(2\mathrm{e}^{-2ax} - \mathrm{e}^{-ax}) .
\end{equation*}
To first order $\mathrm{e}^{-2ax} \approx 1 - 2ax$ and $\mathrm{e}^{-ax}
\approx 1 - ax$, so the bracket is $2 - 4ax - 1 + ax = 1 - 3ax$. At $x =
\pm\SI{0.0345}{\angstrom}$, $3ax = 0.207$: the stiffness is 21\% below its
value at the bottom on the outer side and 21\% above it on the inner.

\solutionto{ex:lo-square}If $X$ has an inverse, so does $X^{\mathsf T}$,
and multiplying the normal equations $X^{\mathsf T}X\vb{w} = X^{\mathsf
T}\vb{y}$ by $(X^{\mathsf T})^{-1}$ on the left gives $X\vb{w} = \vb{y}$.

\solutionto{ex:lo-shrink}As in \cref{sec:lo-ridge}, the component along
$\vb{v}_k$ is $c_k = \vb{v}_k\cdot X^{\mathsf T}\vb{y}/(d_k + \lambda)$
with the penalty and $c_k^0 = \vb{v}_k\cdot X^{\mathsf T}\vb{y}/d_k$
without it, so $c_k = c_k^0\,d_k/(d_k + \lambda)$.

\solutionto{ex:lo-prior}$\sigma_w = \sigma/\sqrt\lambda =
0.2/\sqrt{10^{-3}} = \SI{6.32}{\milli\electronvolt}$; and $\lambda =
\sigma^2/\sigma_w^2 = (0.2/1)^2 = 0.04$.

\solutionto{ex:lo-condition}The mean is $(c/\sigma_b^2)\,b = (1.5/1)\times
2 = 3$, and the variance $\sigma_a^2 - c^2/\sigma_b^2 = 4 - 2.25 = 1.75$.

\solutionto{ex:lo-one}With one reading, $K$ is the number $h^2 +
\sigma^2$ and $\vb{k}_\ast = h^2\mathrm{e}^{-(x - x_1)^2/2\ell^2}$, so
\cref{eq:lo-gp} gives the mean $\vb{k}_\ast y_1/K$, the expression asked
for, and the variance
\begin{equation*}
   h^2 - \frac{h^4\mathrm{e}^{-(x - x_1)^2/\ell^2}}{h^2 + \sigma^2} ,
\end{equation*}
which at $x_1$ is $h^2\sigma^2/(h^2 + \sigma^2)$, below both $h^2$ and
$\sigma^2$, and far away returns to $h^2$.

\solutionto{ex:lo-rate}The rate must be below the stability boundary $2/200 = 0.01$; the best
fixed rate $2/(2 + 200) = 0.0099$; with $\kappa = 100$ the factor is $99/
101 = 0.980$; and the loss falls as its square, so a millionfold takes
$\ln10^6/(-2\ln0.980) = 345$ steps.

\solutionto{ex:lo-kappa}The factor is $(\kappa - 1)/(\kappa + 1)$ and the
loss falls as its square, so $t$ steps bring the factor $\varepsilon$ when
\begin{align*}
   t &= \frac{\ln(1/\varepsilon)}{2\ln[(\kappa + 1)/(\kappa - 1)]}
      && \text{taking logarithms}\\
   &= \frac{\ln(1/\varepsilon)}{2\ln[1 + 2/(\kappa - 1)]}
   \approx \frac{(\kappa - 1)\ln(1/\varepsilon)}{4}
      && \text{writing $(\kappa + 1)/(\kappa - 1) = 1 + 2/(\kappa - 1)$,}
\end{align*}
by $\ln(1 + y) \approx y$ for the small $y = 2/(\kappa - 1)$; for large
$\kappa$ this is $\kappa\ln(1/\varepsilon)/4$.

\solutionto{ex:lo-batch}$\sigma_g/\sqrt b = 0.05$ gives $b = (0.585/0.05)^2
= 137$.

\solutionto{ex:lo-memory}With no gradient $\vb{v}$ is multiplied by $\mu$
each step, so after $t$ steps by $0.9^t = \mathrm{e}^{-1}$ when $t =
-1/\ln0.9 = 9.5$ steps. The friction is $1 - \mu = 0.1$ per step.

\solutionto{ex:lo-polyak}$\sqrt\kappa = 10$, so $\sqrt\mu = 9/11 = 0.818$,
$\mu = 0.669$ and $\eta = 4/(10 + 1)^2 = 0.0331$. Using the asymptotic factors, shrinking by $10^{-4}$
takes approximately $\ln10^{-4}/\ln0.818 = 46$ steps with momentum and $\ln10^{-4}/
\ln0.980 = 461$ without.

\solutionto{ex:lo-adam}$m = (1 - 0.9^{10})\,g = 0.651\,g$ and $s = (1 -
0.999^{10})\,g^2 = 0.00996\,g^2$, so that $\sqrt s$ is
$0.0998\lvert g\rvert$. The corrected step is $\eta g/\lvert g\rvert$; uncorrected it would be
$\eta\times0.651/0.0998 = 6.5\eta$ in size, six and a half times too
large, since $s$ starts far more biased than $m$.

\solutionto{ex:lo-scale}Each $m$ is a weighted sum of the gradients, so
it becomes $cm$; each $s$ a weighted sum of their squares, so $c^2s$, and
$\sqrt{c^2s} = c\sqrt s$ for $c > 0$. The corrections divide both by the
same numbers, and the step $\eta\hat m/\sqrt{\hat s}$ is unchanged.

\solutionto{ex:lo-frames}\SI{2}{\pico\second} is 200 frames, about $200/
109 = 1.8$ independent ones; \SI{20}{\pico\second}, about 18.

\solutionto{ex:lo-neighbour}The frames on either side differ from it by
\SI{10}{\femto\second} of motion, much less than the memory of the run,
so their features and energies are nearly its own. A model that fits its
training frames and interpolates smoothly between their features can
therefore predict this validation frame with a similarly small miss.
A new run need not supply such close neighbours. The random split then
tests local interpolation more strongly than transfer to new configurations.

\solutionto{ex:lo-plan}Split by whole runs, never by frames: for
instance, if there are five runs at each of the ten temperatures, put
three at each temperature in training, one in validation and one in test.
This gives 30, 10 and 10 runs, with every temperature represented in each set. No two sets may share a run,
and the scaling of the features, the choice of settings and the early
stop come from the training and validation sets alone; the test runs are
used once, at the end.
